\documentclass[11pt]{article}
% Préambule commun aux cours de la formation C++.
% Inclus par chaque cours.tex via % Préambule commun aux cours de la formation C++.
% Inclus par chaque cours.tex via % Préambule commun aux cours de la formation C++.
% Inclus par chaque cours.tex via \input{preambule}.

\usepackage[utf8]{inputenc}
\usepackage[T1]{fontenc}
\usepackage[french]{babel}
\usepackage{lmodern}
\usepackage{microtype}

\usepackage{geometry}
\geometry{a4paper, margin=2.5cm}

\usepackage{xcolor}
\usepackage{amsmath, amssymb}
\usepackage{enumitem}
\usepackage{hyperref}
\hypersetup{
    colorlinks=true,
    linkcolor=blue!60!black,
    urlcolor=blue!60!black,
    citecolor=blue!60!black,
}

% --- Coloration syntaxique du C++ via listings -----------------------------
\usepackage{listings}

\definecolor{codebg}{rgb}{0.97,0.97,0.95}
\definecolor{codekw}{rgb}{0.10,0.10,0.60}
\definecolor{codecomment}{rgb}{0.20,0.50,0.20}
\definecolor{codestring}{rgb}{0.65,0.15,0.15}
\definecolor{coderule}{rgb}{0.80,0.80,0.80}

% Permet les caractères accentués (UTF-8) à l'intérieur des blocs de code.
\lstset{
    literate=%
        {é}{{\'e}}1 {è}{{\`e}}1 {ê}{{\^e}}1 {ë}{{\"e}}1
        {à}{{\`a}}1 {â}{{\^a}}1 {ä}{{\"a}}1
        {î}{{\^i}}1 {ï}{{\"i}}1
        {ô}{{\^o}}1 {ö}{{\"o}}1
        {ù}{{\`u}}1 {û}{{\^u}}1 {ü}{{\"u}}1
        {ç}{{\c c}}1 {É}{{\'E}}1 {È}{{\`E}}1 {À}{{\`A}}1
}

\lstdefinestyle{cpp}{
    language=C++,
    backgroundcolor=\color{codebg},
    basicstyle=\ttfamily\small,
    keywordstyle=\color{codekw}\bfseries,
    commentstyle=\color{codecomment}\itshape,
    stringstyle=\color{codestring},
    showstringspaces=false,
    numbers=left,
    numberstyle=\tiny\color{gray},
    numbersep=8pt,
    frame=single,
    rulecolor=\color{coderule},
    breaklines=true,
    tabsize=4,
    morekeywords={constexpr,noexcept,override,final,nullptr,auto,
        std,string,vector,size_t,uint64_t,int64_t,optional,namespace}
}
\lstset{style=cpp}

% Raccourci pour du code en ligne, en police machine à écrire.
\newcommand{\code}[1]{\texttt{#1}}

% Boîte « À retenir »
\usepackage{tcolorbox}
\newtcolorbox{aretenir}[1][]{
    colback=yellow!8, colframe=orange!70!black,
    title=#1, fonttitle=\bfseries, sharp corners, boxrule=0.5pt
}
\newtcolorbox{piege}[1][Piège]{
    colback=red!5, colframe=red!60!black,
    title=#1, fonttitle=\bfseries, sharp corners, boxrule=0.5pt
}
.

\usepackage[utf8]{inputenc}
\usepackage[T1]{fontenc}
\usepackage[french]{babel}
\usepackage{lmodern}
\usepackage{microtype}

\usepackage{geometry}
\geometry{a4paper, margin=2.5cm}

\usepackage{xcolor}
\usepackage{amsmath, amssymb}
\usepackage{enumitem}
\usepackage{hyperref}
\hypersetup{
    colorlinks=true,
    linkcolor=blue!60!black,
    urlcolor=blue!60!black,
    citecolor=blue!60!black,
}

% --- Coloration syntaxique du C++ via listings -----------------------------
\usepackage{listings}

\definecolor{codebg}{rgb}{0.97,0.97,0.95}
\definecolor{codekw}{rgb}{0.10,0.10,0.60}
\definecolor{codecomment}{rgb}{0.20,0.50,0.20}
\definecolor{codestring}{rgb}{0.65,0.15,0.15}
\definecolor{coderule}{rgb}{0.80,0.80,0.80}

% Permet les caractères accentués (UTF-8) à l'intérieur des blocs de code.
\lstset{
    literate=%
        {é}{{\'e}}1 {è}{{\`e}}1 {ê}{{\^e}}1 {ë}{{\"e}}1
        {à}{{\`a}}1 {â}{{\^a}}1 {ä}{{\"a}}1
        {î}{{\^i}}1 {ï}{{\"i}}1
        {ô}{{\^o}}1 {ö}{{\"o}}1
        {ù}{{\`u}}1 {û}{{\^u}}1 {ü}{{\"u}}1
        {ç}{{\c c}}1 {É}{{\'E}}1 {È}{{\`E}}1 {À}{{\`A}}1
}

\lstdefinestyle{cpp}{
    language=C++,
    backgroundcolor=\color{codebg},
    basicstyle=\ttfamily\small,
    keywordstyle=\color{codekw}\bfseries,
    commentstyle=\color{codecomment}\itshape,
    stringstyle=\color{codestring},
    showstringspaces=false,
    numbers=left,
    numberstyle=\tiny\color{gray},
    numbersep=8pt,
    frame=single,
    rulecolor=\color{coderule},
    breaklines=true,
    tabsize=4,
    morekeywords={constexpr,noexcept,override,final,nullptr,auto,
        std,string,vector,size_t,uint64_t,int64_t,optional,namespace}
}
\lstset{style=cpp}

% Raccourci pour du code en ligne, en police machine à écrire.
\newcommand{\code}[1]{\texttt{#1}}

% Boîte « À retenir »
\usepackage{tcolorbox}
\newtcolorbox{aretenir}[1][]{
    colback=yellow!8, colframe=orange!70!black,
    title=#1, fonttitle=\bfseries, sharp corners, boxrule=0.5pt
}
\newtcolorbox{piege}[1][Piège]{
    colback=red!5, colframe=red!60!black,
    title=#1, fonttitle=\bfseries, sharp corners, boxrule=0.5pt
}
.

\usepackage[utf8]{inputenc}
\usepackage[T1]{fontenc}
\usepackage[french]{babel}
\usepackage{lmodern}
\usepackage{microtype}

\usepackage{geometry}
\geometry{a4paper, margin=2.5cm}

\usepackage{xcolor}
\usepackage{amsmath, amssymb}
\usepackage{enumitem}
\usepackage{hyperref}
\hypersetup{
    colorlinks=true,
    linkcolor=blue!60!black,
    urlcolor=blue!60!black,
    citecolor=blue!60!black,
}

% --- Coloration syntaxique du C++ via listings -----------------------------
\usepackage{listings}

\definecolor{codebg}{rgb}{0.97,0.97,0.95}
\definecolor{codekw}{rgb}{0.10,0.10,0.60}
\definecolor{codecomment}{rgb}{0.20,0.50,0.20}
\definecolor{codestring}{rgb}{0.65,0.15,0.15}
\definecolor{coderule}{rgb}{0.80,0.80,0.80}

% Permet les caractères accentués (UTF-8) à l'intérieur des blocs de code.
\lstset{
    literate=%
        {é}{{\'e}}1 {è}{{\`e}}1 {ê}{{\^e}}1 {ë}{{\"e}}1
        {à}{{\`a}}1 {â}{{\^a}}1 {ä}{{\"a}}1
        {î}{{\^i}}1 {ï}{{\"i}}1
        {ô}{{\^o}}1 {ö}{{\"o}}1
        {ù}{{\`u}}1 {û}{{\^u}}1 {ü}{{\"u}}1
        {ç}{{\c c}}1 {É}{{\'E}}1 {È}{{\`E}}1 {À}{{\`A}}1
}

\lstdefinestyle{cpp}{
    language=C++,
    backgroundcolor=\color{codebg},
    basicstyle=\ttfamily\small,
    keywordstyle=\color{codekw}\bfseries,
    commentstyle=\color{codecomment}\itshape,
    stringstyle=\color{codestring},
    showstringspaces=false,
    numbers=left,
    numberstyle=\tiny\color{gray},
    numbersep=8pt,
    frame=single,
    rulecolor=\color{coderule},
    breaklines=true,
    tabsize=4,
    morekeywords={constexpr,noexcept,override,final,nullptr,auto,
        std,string,vector,size_t,uint64_t,int64_t,optional,namespace}
}
\lstset{style=cpp}

% Raccourci pour du code en ligne, en police machine à écrire.
\newcommand{\code}[1]{\texttt{#1}}

% Boîte « À retenir »
\usepackage{tcolorbox}
\newtcolorbox{aretenir}[1][]{
    colback=yellow!8, colframe=orange!70!black,
    title=#1, fonttitle=\bfseries, sharp corners, boxrule=0.5pt
}
\newtcolorbox{piege}[1][Piège]{
    colback=red!5, colframe=red!60!black,
    title=#1, fonttitle=\bfseries, sharp corners, boxrule=0.5pt
}


\title{\textbf{Chapitre 1 — Fondamentaux : compiler, typer, choisir}\\[0.3em]
       \large Le cœur du langage, sans STL avancée}
\author{Formation C++ moderne}
\date{}

\begin{document}
\maketitle

\begin{abstract}
\noindent Ce premier chapitre réactive tes réflexes C++ et installe les bonnes
habitudes qu'on gardera toute la formation. Tu connais déjà la programmation :
on ne réexplique pas ce qu'est une boucle, on se concentre sur \emph{le cœur du
langage} et sur ses idiomes \emph{modernes} (C++17/20). À la fin tu sauras
compiler proprement avec les warnings, choisir tes types sans te faire piéger
par l'overflow, distinguer copie et référence, écrire et surcharger des
fonctions, manier \code{const}, structurer ton flot de contrôle, et signaler
une erreur sans crasher. On reste \emph{volontairement} sur le langage :
\code{std::string}, \code{std::vector}, l'organisation \code{.hpp}/\code{.cpp},
les namespaces et le RAII arrivent au chapitre~2. Le tout testé avec un petit
harnais maison, \code{check.hpp}.
\end{abstract}

\tableofcontents
\bigskip\hrule\bigskip

% ===========================================================================
\section{Compilation, flags et warnings}

Un programme C++ passe par trois étapes avant de devenir exécutable :

\begin{enumerate}[noitemsep]
    \item \textbf{Préprocesseur} : déroule les \code{\#include}, les macros.
    \item \textbf{Compilation} : chaque \code{.cpp} (\emph{unité de traduction})
          devient un fichier objet \code{.o}.
    \item \textbf{Édition de liens} (\emph{linker}) : assemble les \code{.o} et
          les bibliothèques en un exécutable.
\end{enumerate}

La ligne de compilation qu'on utilisera tout le chapitre :

\begin{lstlisting}
g++ -std=c++20 -Wall -Wextra -Wpedantic mon_prog.cpp -o mon_prog
\end{lstlisting}

\begin{itemize}[noitemsep]
    \item \code{-std=c++20} : on travaille en C++ moderne.
    \item \code{-Wall -Wextra -Wpedantic} : active un maximum d'avertissements.
\end{itemize}

\begin{aretenir}[Lis tes warnings]
Un warning n'est pas du bruit : c'est le compilateur qui te dit « ceci est
probablement un bug ». Conversion qui perd de l'information, variable non
initialisée, comparaison signé/non-signé\ldots\ Traite chaque warning comme une
erreur à comprendre. On apprend énormément en les lisant.
\end{aretenir}

% ===========================================================================
\section{Structure d'un programme}

\begin{lstlisting}
#include <iostream>

int main() {
    std::cout << "Bonjour\n";
    return 0;   // 0 = succes (convention Unix). Non-zero = erreur.
}
\end{lstlisting}

\code{main} est le point d'entrée. Son entier de retour est le \emph{code de
sortie} vu par le shell (\code{echo \$?}). \code{return 0} est implicite dans
\code{main} (mais explicite ailleurs). On exploitera ce code de retour : nos
fichiers de test feront \code{return check::report();} (0 si tout passe, 1
sinon). C'est exactement ce que vérifie le shell pour savoir si ta batterie de
tests est passée.

% ===========================================================================
\section{Entrées/sorties : \code{cout}, \code{cin}, \code{cerr}}

Les flux standard vivent dans \code{<iostream>}. \code{std::cout} est la sortie
standard, \code{std::cerr} la sortie d'erreur (séparée, non tamponnée),
\code{std::cin} l'entrée standard. On y envoie avec \code{<<} et on lit avec
\code{>>}.

\begin{lstlisting}
std::cout << "x = " << x << '\n';   // sortie standard
std::cerr << "Erreur : ...\n";       // sortie d'erreur (separee)
int n;
std::cin >> n;                        // lecture (verifie l'etat du flux !)
\end{lstlisting}

\begin{piege}[\code{'\textbackslash n'} vs \code{std::endl}]
\code{std::endl} insère un retour à la ligne \emph{et force un flush} du tampon
de sortie (opération coûteuse). Pour un simple saut de ligne, préfère
\code{'\textbackslash n'}. Garde \code{std::endl} pour quand tu veux vraiment
forcer l'écriture immédiate.
\end{piege}

Lecture robuste : \code{std::cin >> n} peut \emph{échouer} (entrée non
numérique). L'expression d'extraction renvoie le flux, convertible en
\code{bool} : teste-la directement.

\begin{lstlisting}
int n;
if (std::cin >> n) {
    // lecture reussie : n contient la valeur
} else {
    std::cerr << "entree invalide\n";
}
\end{lstlisting}

% ===========================================================================
\section{Types fondamentaux, overflow, \code{static\_cast} et \code{auto}}

Entiers : \code{int} (au moins 16 bits, en pratique 32), \code{long long} (au
moins 64), versions \code{unsigned}. Flottants : \code{double} par défaut.
\code{bool}, \code{char}.

\begin{piege}[Débordement d'entier (overflow)]
Un calcul qui dépasse la capacité du type « boucle » silencieusement vers une
valeur fausse (souvent négative), \emph{sans} erreur à l'exécution. Exemple :
\code{850000 * 850000} déborde un \code{int} 32 bits ($\approx 2{,}1$
milliards). La parade : un type assez grand (\code{long long}), et attention,
\code{int * int} se calcule \emph{en} \code{int} même si tu ranges le résultat
dans un \code{long long} — il faut qu'un opérande soit déjà 64 bits.
\end{piege}

Littéraux pratiques : séparateur \code{4'000'000}, suffixes \code{42LL}
(\code{long long}), \code{42u} (\code{unsigned}).

\paragraph{Convertir un type : \code{static\_cast}.} Le suffixe \code{LL} ne
marche que sur un \emph{littéral} écrit en dur (\code{42LL}). Pour convertir une
\emph{variable} d'un type vers un autre, on utilise un \textbf{cast} explicite :

\begin{lstlisting}
int n = 100000;
long long produit = static_cast<long long>(n) * n;  // calcul fait en 64 bits
\end{lstlisting}

\code{static\_cast<TypeCible>(expression)} convertit \code{expression} vers
\code{TypeCible}. Ici, dès que \emph{l'un} des opérandes est un \code{long long},
toute la multiplication se fait en 64 bits — plus de débordement.

\begin{aretenir}[\code{static\_cast} plutôt que le cast à la C]
Tu verras parfois \code{(long long)n} (style hérité du C). En C++ moderne,
préfère \code{static\_cast<long long>(n)} : plus visible, plus facile à
rechercher, et le compilateur vérifie que la conversion a un sens.
\end{aretenir}

\paragraph{Le typage est statique.} Chaque variable a un type fixé à la
\emph{compilation} et ne change jamais. \code{int n = 3;} fait de \code{n} un
entier pour toujours : tu ne peux pas lui réaffecter du texte plus tard. Le
compilateur connaît donc tous les types avant même d'exécuter le programme,
ce qui lui permet d'attraper des erreurs très tôt et de générer du code rapide.

\paragraph{Déclarer = donner un type.} On déclare une variable en écrivant son
type devant son nom : \code{double moyenne = 0.0;}. Tant qu'une variable n'est
pas déclarée, elle n'existe pas — pas de création « à la volée » au milieu d'un
calcul, contrairement à Python.

\code{auto} demande au compilateur de \emph{déduire} le type à partir de la
valeur d'initialisation. La variable reste fortement typée ; tu t'épargnes
juste l'écriture du type :

\begin{lstlisting}
auto i = 42;              // i est un int (deduit du litteral 42)
auto x = 3.14;            // x est un double
auto& r = tableau[0];     // REFERENCE vers l'element (le & compte !)
\end{lstlisting}

\begin{piege}[\code{auto} sans \code{\&} fait une copie]
\code{auto x = v[0];} crée une \emph{copie} de l'élément : modifier \code{x} ne
touche pas \code{v}. \code{auto\& r = v[0];} crée une \emph{référence} : modifier
\code{r} modifie \code{v[0]}. Cette nuance — un seul caractère — est l'une des
plus importantes du langage. On y revient en détail à la section sur les
références.
\end{piege}

\begin{aretenir}[Quand utiliser \code{auto} ?]
Utilise-le quand le type est évident ou verbeux. Évite-le quand nommer le type
\emph{clarifie} l'intention. Et décide toujours consciemment : veux-tu une copie
(\code{auto}) ou une référence (\code{auto\&}, \code{const auto\&}) ?
\end{aretenir}

% ===========================================================================
\section{\code{const} et const-correctness}

\code{const} dit « ceci ne changera pas ». C'est une garantie pour le lecteur
\emph{et} pour le compilateur.

\begin{lstlisting}
const int n = 10;          // valeur figee : n = 11; serait refuse a la compilation
const double pi = 3.14159;
\end{lstlisting}

Sur un paramètre de fonction, \code{const} documente que l'argument ne sera pas
modifié. Combiné à la référence (\code{const T\&}, section suivante), il permet
de \emph{lire} un objet sans le copier ni le modifier :

\begin{lstlisting}
double aire_disque(const double rayon);     // promesse : rayon ne change pas
\end{lstlisting}

\begin{aretenir}[Const-correctness]
Marque \code{const} tout ce qui ne doit pas changer : variables locales,
paramètres en lecture seule, et plus tard les méthodes qui ne modifient pas
l'objet. C'est plus sûr (le compilateur refuse les modifications accidentelles)
et plus expressif. Avant-goût : \code{constexpr} (calcul à la \emph{compilation})
qu'on verra plus tard.
\end{aretenir}

% ===========================================================================
\section{Valeur ou référence : LA distinction à maîtriser}

C'est sans doute le point le plus important du chapitre. Prends ton temps ici.

\subsection*{Par défaut, C++ copie}

Quand tu écris \code{b = a} pour deux variables, ou quand tu passes une variable
à une fonction, C++ fait par défaut une \textbf{copie} : \code{b} est un
\emph{nouvel} objet, indépendant de \code{a}.

\begin{lstlisting}
int a = 10;
int b = a;     // b est une COPIE : b vaut 10, mais c'est un autre entier
b = 99;        // on change b...
// a vaut toujours 10 -- a et b sont independants
\end{lstlisting}

C'est l'inverse du réflexe Python, où une affectation partage une référence vers
le même objet. En C++, l'affectation par défaut, c'est de la copie.

\subsection*{Une référence ne copie pas : c'est un autre nom}

Une \textbf{référence} \code{T\&} n'est pas un nouvel objet : c'est un
\emph{alias}, un second nom pour une variable qui existe déjà. Modifier la
référence, c'est modifier la variable d'origine. Une référence doit être liée à
sa création et ne peut jamais être re-liée vers autre chose ensuite.

\begin{lstlisting}
int a = 10;
int& r = a;    // r est un AUTRE NOM pour a (pas une copie)
r = 99;        // on ecrit a travers r...
// a vaut maintenant 99 ! r et a designent le meme entier
\end{lstlisting}

\begin{aretenir}[La règle mentale]
\code{int b = a;} → « fais-moi un \emph{nouveau} truc égal à \code{a} » (copie).\\
\code{int\& r = a;} → « \code{r} \emph{est} \code{a}, juste un autre nom »
(référence).
\end{aretenir}

\subsection*{À quoi ça sert : modifier un argument}

Une fonction reçoit ses arguments par copie par défaut, donc elle ne peut pas
modifier les variables de l'appelant. Pour qu'elle le puisse, on passe par
référence :

\begin{lstlisting}
void incremente_copie(int x) { x += 1; }    // modifie une COPIE locale, sans effet
void incremente_ref(int& x)  { x += 1; }    // modifie l'original

int n = 5;
incremente_copie(n);   // n vaut toujours 5
incremente_ref(n);     // n vaut maintenant 6
\end{lstlisting}

L'échange de deux variables illustre parfaitement le besoin de référence :

\begin{lstlisting}
void swap_ints(int& a, int& b) {   // SANS le &, on echangerait des copies : inutile
    int t = a;
    a = b;
    b = t;
}
\end{lstlisting}

\subsection*{À quoi ça sert aussi : éviter une copie coûteuse}

Copier un \code{int} est gratuit. Copier un gros objet (une longue chaîne, une
grande liste de valeurs), non. Si une fonction veut juste \emph{lire} un gros
objet, on le passe par \code{const T\&} : pas de copie, et le \code{const}
interdit toute modification accidentelle.

\begin{lstlisting}
std::size_t longueur(const std::string& s);   // ni copie, ni modification
\end{lstlisting}

\begin{aretenir}[Trois façons de passer un argument]
\code{T} (valeur) : la fonction reçoit une copie. Pour les petits types ou
quand on veut vraiment une copie.\\
\code{T\&} (référence) : la fonction peut modifier l'original.\\
\code{const T\&} (référence constante) : la fonction lit l'original sans le
copier ni le modifier. \textbf{Le choix par défaut pour les gros objets en
lecture seule.}
\end{aretenir}

\begin{aretenir}[Avant-goût : \code{std::string} et \code{std::vector}]
Quelques exercices de ce chapitre frôlent deux types très utiles, qu'on détaille
au chapitre~2. Pour l'instant, retiens l'essentiel :
\begin{itemize}[noitemsep]
    \item \code{std::string} est une chaîne de caractères qui gère sa propre
          taille : \code{s.size()}, concaténation avec \code{+}, comparaison
          avec \code{==}.
    \item \code{std::vector<T>} est une liste qui grandit toute seule :
          \code{v.push\_back(x)} ajoute en fin, \code{v[i]} accède au
          $i$-ème élément, \code{v.size()} donne le nombre d'éléments.
\end{itemize}
Ce sont de « gros objets » : passe-les par \code{const T\&} en lecture. Le reste
(mémoire, capacité, parcours) attend le chapitre~2.
\end{aretenir}

\subsection*{Et les pointeurs ?}

Les \textbf{pointeurs} \code{T*} (qui peuvent valoir \code{nullptr} et qu'on
déréférence avec \code{*}) sont l'outil de plus bas niveau, approfondi dans un
chapitre dédié plus loin dans la formation. Pour l'instant : \emph{quand une
référence suffit, préfère la référence} — elle est plus sûre (jamais nulle,
jamais re-liée par erreur).

% ===========================================================================
\section{Fonctions : déclaration, surcharge, passage des paramètres}

Une fonction se \emph{déclare} (signature : type de retour, nom, paramètres)
puis se \emph{définit} (corps). On peut déclarer avant d'utiliser et définir
plus loin.

\begin{lstlisting}
int carre(int n);            // declaration (prototype)

int main() {
    return carre(7);         // on peut appeler : la declaration suffit
}

int carre(int n) { return n * n; }   // definition
\end{lstlisting}

\textbf{Surcharge} : plusieurs fonctions de même nom, différenciées par le
\emph{type} ou le \emph{nombre} de leurs paramètres. Le compilateur choisit la
bonne version selon les arguments de l'appel (\emph{overload resolution}).

\begin{lstlisting}
std::string describe(int n);
std::string describe(double d);
std::string describe(bool b);
std::string describe(const std::string& s);

describe(42);        // appelle la version int
describe(3.14);      // appelle la version double
describe("salut");   // appelle la version const std::string&
\end{lstlisting}

\begin{piege}[Surcharge ambiguë]
Si deux surcharges sont aussi « bonnes » l'une que l'autre pour un appel donné,
le compilateur refuse : \emph{call of overloaded \ldots is ambiguous}. Méfie-toi
des conversions implicites (un \code{char} qui pourrait aller vers \code{int}
\emph{ou} \code{double}). Des types de paramètres bien distincts évitent
l'ambiguïté.
\end{piege}

\textbf{Passage des paramètres} — le choix qui compte :

\begin{center}
\begin{tabular}{lll}
\hline
\textbf{Forme} & \textbf{Quand} & \textbf{Effet} \\
\hline
\code{T} (valeur) & petits types, copie voulue & copie l'argument \\
\code{const T\&} & gros objets en lecture seule & pas de copie, non modifiable \\
\code{T\&} & on veut \emph{modifier} l'argument & pas de copie, modifiable \\
\hline
\end{tabular}
\end{center}

\begin{aretenir}[Règle par défaut]
Passe les types fondamentaux (\code{int}, \code{double}\ldots) \emph{par
valeur}. Passe les gros objets \emph{par \code{const T\&}} si tu les lis, par
\code{T\&} si tu les modifies. Copier un gros objet sans raison, c'est du
gâchis.
\end{aretenir}

% ===========================================================================
\section{Structures de contrôle}

\code{if}/\code{else}, \code{switch}, \code{while}, \code{do-while},
\code{for}, \code{break}, \code{continue} : rien de surprenant pour un
programmeur. Deux nouveautés modernes méritent ton attention.

\paragraph{Init-statement (C++17).} \code{if} et \code{switch} acceptent une
instruction d'initialisation avant la condition, séparée par \code{;}. La
variable déclarée n'existe que dans le \code{if}/\code{else} : portée minimale,
moins de fuites de noms.

\begin{lstlisting}
if (int q = a / b; q > 0) {   // q n'existe que dans ce if/else
    utiliser(q);
} else {
    std::cerr << "quotient nul ou negatif\n";
}
\end{lstlisting}

\paragraph{\code{switch} et \code{[[fallthrough]]}.} Dans un \code{switch}, un
\code{case} sans \code{break} « tombe » dans le suivant. C'est souvent un bug ;
\code{-Wextra} t'avertit. Quand la chute est \emph{voulue}, documente-la avec
l'attribut \code{[[fallthrough]]} (C++17) : l'intention est claire et le warning
disparaît.

\begin{lstlisting}
switch (cmd) {
    case 'a': /* ... */ break;
    case 'b': [[fallthrough]];   // chute volontaire, documentee
    case 'c': /* ... */ break;
    default:  /* ... */ break;
}
\end{lstlisting}

\begin{piege}[\code{switch} et déclaration de variable]
Déclarer une variable initialisée directement dans un \code{case} sans bloc
\code{\{ \}} provoque une erreur (\emph{crosses initialization}). Entoure le
corps du \code{case} d'accolades \code{\{ \ldots \}} pour lui donner sa propre
portée.
\end{piege}

% ===========================================================================
\section{Portée et durée de vie}

Une variable locale \textbf{naît} à sa déclaration et \textbf{meurt} à la fin de
son bloc — l'accolade fermante \code{\}}. C'est automatique : tu n'as rien à
libérer toi-même.

\begin{lstlisting}
void f() {
    int a = 1;            // a nait ici
    {
        int b = 2;        // b nait ici
        std::cout << a + b;
    }                     // b MEURT ici (fin de son bloc)
    // b n'existe plus, a existe encore
}                         // a MEURT ici
\end{lstlisting}

\paragraph{Pile et tas (survol).} Les variables locales comme \code{a} et
\code{b} vivent sur la \textbf{pile} (\emph{stack}) : leur création et leur
destruction sont gérées automatiquement à l'entrée et à la sortie du bloc, c'est
quasi gratuit. Certains objets demandent en plus de la mémoire sur le
\textbf{tas} (\emph{heap}) pour stocker leurs données — mais, lorsqu'ils sont
bien conçus, ils s'en occupent eux-mêmes : ils l'acquièrent à leur création et
la libèrent à leur destruction. On entre dans le détail (pile vs tas, allocation
dynamique, le mécanisme RAII) au chapitre~2 et au chapitre dédié aux pointeurs.

\paragraph{Conseil pratique.} Déclare chaque variable au plus près de son usage,
dans le bloc le plus restreint possible. Une variable qui n'existe que là où on
s'en sert, c'est moins d'erreurs et un code plus lisible. L'init-statement de la
section précédente sert exactement ça.

\begin{piege}[Shadowing]
Une variable interne peut en « masquer » une externe du même nom : à l'intérieur
du bloc, c'est la plus interne qui compte, et l'externe redevient visible à la
sortie. Le warning \code{-Wshadow} t'avertit. Évite : ça prête à confusion.
\end{piege}

% ===========================================================================
\section{Gestion d'erreurs : codes de retour vs exceptions}

Deux grandes approches pour signaler qu'une opération a échoué.

\paragraph{1. Codes de retour.} La fonction renvoie un \code{bool}, un code, ou
écrit dans un paramètre de sortie passé par référence. Léger et explicite, mais
\emph{oblige l'appelant à tester à chaque appel} — un oubli passe inaperçu.

\begin{lstlisting}
bool diviser(int a, int b, int& resultat) {
    if (b == 0) return false;       // echec signale par le code de retour
    resultat = a / b;
    return true;
}
\end{lstlisting}

\paragraph{2. Exceptions.} Pour les cas \emph{exceptionnels}, on \code{throw}
une exception ; le flot normal est interrompu et l'erreur remonte la pile
d'appels jusqu'à un \code{try}/\code{catch} qui la gère. Les types d'exception
standard vivent dans \code{<stdexcept>} : \code{std::invalid\_argument},
\code{std::domain\_error}, \code{std::out\_of\_range}\ldots\ Ils prennent un
message, qu'on relit avec \code{e.what()}.

\begin{lstlisting}
#include <stdexcept>

double racine(double x) {
    if (x < 0) throw std::domain_error("racine d'un nombre negatif");
    // ... calcul ...
}

// appel protege :
try {
    double r = racine(valeur);
    std::cout << r << '\n';
} catch (const std::domain_error& e) {   // on attrape par const reference
    std::cerr << e.what() << '\n';
}
\end{lstlisting}

\begin{aretenir}[Codes ou exceptions ?]
Réserve les exceptions aux situations \emph{anormales} (entrée invalide,
précondition violée) — pas au flot de contrôle ordinaire. Pour un échec attendu
et fréquent, un code de retour (ou plus tard \code{std::optional<T>}, « une
valeur \emph{ou} rien », sans sentinelle \code{-1}) est souvent plus clair.
Attrape toujours une exception \emph{par \code{const} référence}
(\code{catch (const std::exception\& e)}).
\end{aretenir}

% ===========================================================================
\section{Nommage et style}

\begin{itemize}[noitemsep]
    \item Sois cohérent : \code{snake\_case} pour les fonctions/variables est un
          choix répandu (le nôtre) ; \code{CamelCase} pour les types.
    \item Des noms qui \emph{disent l'intention} : \code{prev}/\code{current}
          plutôt que \code{cpt1}/\code{cpt2}.
    \item Indentation 4 espaces, accolades cohérentes.
    \item Une variable par responsabilité ; \code{const} dès que possible.
    \item Un commentaire explique le \emph{pourquoi}, pas le \emph{quoi} évident.
\end{itemize}

\begin{aretenir}[Le fil rouge du chapitre]
Compile avec les warnings et corrige-les. Choisis tes types en pensant à
l'overflow (\code{static\_cast} vers \code{long long} quand il faut). Distingue
copie et référence, et passe les gros objets par \code{const\&}. Structure
proprement ton flot de contrôle. Signale les erreurs avec un code de retour ou
une exception, selon que l'échec est ordinaire ou exceptionnel. Ces réflexes te
suivront tout le reste de la formation.
\end{aretenir}

\bigskip\hrule\bigskip
\noindent\textbf{À toi de jouer.} Les exercices du dossier \code{exercices/}
suivent l'ordre de ce chapitre, en difficulté croissante :
\code{01-echauffement} (premières fonctions, types et code retour),
\code{02-temperatures} (calcul, conversions, choix de type),
\code{03-valeur-reference} (copie vs \code{T\&} : modifier ou non l'argument de
l'appelant), \code{04-const-exceptions} (\code{const\&} pour lire, \code{\&}
pour modifier, et \code{throw} pour signaler l'erreur), \code{05-surcharge}
(une fonction \code{describe}, plusieurs surcharges choisies par le
compilateur), \code{06-fizzbuzz} (structures de contrôle), puis deux exercices
plus autonomes : un \emph{dispatch} sur commande (\code{switch}, init-statement,
\code{[[fallthrough]]}) et une \emph{racine sûre} (gérer le cas négatif par
exception ou code de retour). Chaque énoncé introduit, au besoin, l'outil dont
tu auras besoin. Le mini-projet de synthèse — une \emph{boîte à outils
numériques} — assemble fonctions, choix de types, \code{const}-correctness et
gestion d'erreurs dans un petit programme complet, \code{main} compris. Bon
courage !

\end{document}
